\chapter{Bloques}

Sea $k$ la cantidad de tópicos y \(n\) la cantidad de documentos del corpus.

La coherencia de un tópico se calcula como:
\begin{equation}
  C(t) = \sum_{i<j} \log \frac{D(w_i, w_j) + 1}{D(w_j)}
  \label{eq:coherencia}
\end{equation}

El modelo se descompone en dos pasos:
\begin{align}
  \theta_d &\sim \mathrm{Dir}(\alpha) \\
  z_{d,n} &\sim \mathrm{Mult}(\theta_d) \nonumber
\end{align}

En forma resumida:
\[
  p(w \mid d) = \sum_{k} p(w \mid k) \, p(k \mid d)
\]

\begin{figure}[htbp]
  \centering
  \includegraphics[width=0.8\textwidth]{figs/ejemplo.png}
  \caption{Distribución de \textbf{tópicos}}
  \label{fig:distribucion}
\end{figure}

\begin{figure}[htbp]
  \centering
  \includegraphics[width=0.4\textwidth]{figs/ejemplo.png}
  \includegraphics[width=0.4\textwidth]{figs/ejemplo.png}
\end{figure}

\begin{table}[htbp]
  \centering
  \caption{Tamaño del corpus}
  \label{tab:corpus}
  \begin{tabular}{ll}
    \toprule
    Sitio & Documentos \\
    \midrule
    Foro & 1200 \\
    \bottomrule
  \end{tabular}
\end{table}

\begin{lstlisting}[language=Python]
precio = "$10"  # no es \textbf{formato}
print(precio)
\end{lstlisting}

En la figura~\ref{fig:distribucion}, la tabla~\ref{tab:corpus} y la ecuación~\ref{eq:coherencia} se resume lo anterior.
